\subsection{Sqrt decomposition}

Ideia principal: dividir em buckets de tamanho raiz.

Bota tamanho const pro tamanho de bloco.

\subsubsection{Algoritmo de Mo}

Responder $Q$ queries de maneira offline em $O((N+Q) \sqrt{N})$

Parece pior, mas a questão é que se fosse fazer do jeito normal, o passo
de merge teria complicações.

Aqui o código pro DQUERY do SPOJ. Cada query pede o número de valores distintos
num dado intervalo.

\inccode{estrut_dados/dquery.cpp}

\subsubsection{Mo com update}

Cada query além do l, r e id, tem o t, que indica quantos updates tem antes dela

Ordenam-se as queries por bloco do l. Empates são resolvidos pelo bloco do r.
Empates são resolvidos pelo t

Blocos de tamanho  $(2n^2)^{1/3}$

Depois disso faz:

\inccode{estrut_dados/small_mo_update.cpp}

em que o viagem no tempo percorre as updates no intervalo 
de tempo de q.t a t, oq pode cancelar algumas updates ou ativar algumas updates

Segue o código de uma questão que pedia, num intervalo, a soma dos caras distintos

\inccode{estrut_dados/XXXSPOJ_sum_of_distinct.cpp}
